% !TEX root = ../main.tex
% =============================================================================
% CHAPTER 6 - DISCUSSION                                          (target ~6-7 pp)
%
% VOICE : impersonal passive, matching Chapters 4 and 5. Never "we", never "I".
% TENSE : present for interpretation and claims; simple past when referring back
%         to what was found ("as Section 5.3 showed").
%
% This chapter synthesises across studies. Any sentence that could have been
% written from a single table belongs in Chapter 5.
% =============================================================================

\chapter{Discussion}\label{ch:discussion}

\section{Synthesis of Findings}\label{sec:d_synthesis}

Taken together, the experiments reveal three broader findings. First, self-supervised pretraining produced resting-state EEG representations that supported participant level clinical classification, but their usefulness was task dependent. The proposed BrainLM-EEG variants were competitive across several tasks, with their clearest relative advantage observed on FTD under LOSO. Strong performance on AD across multiple model families further suggests that clinically relevant information in that dataset was accessible to a range of modelling approaches.

Second, within-dataset and cross-dataset generalisation were distinct challenges. Representations that transferred successfully to unseen participants within the same dataset did not retain the same level of performance on the independently acquired SCZ dataset. This indicates that participant level evaluation alone provides limited evidence of robustness to changes in cohort and acquisition conditions.

Third, the usefulness of the learned representations depended on design and extraction choices. Masking configuration and encoder depth affected downstream performance, while the CLS-token and reconstruction analyses showed that information relevant to reconstruction and classification was retained beyond a single final representation. The geometry-aware and microstate-informed variants produced improvements in individual settings, but neither provided a consistent advantage across tasks.

Overall, self-supervised pretraining produced reusable EEG representations that performed competitively across multiple clinical tasks and showed that their downstream utility was influenced by the clinical task, evaluation protocol, representation design, and degree of dataset shift. Consistent superiority over supervised or classical alternatives was not observed.


\section{Task Dependence of the Learned Representations}\label{sec:d_clinical}

A central finding is that the amount of participant level separation differed substantially across the four downstream tasks, while differences among model families were generally smaller and less consistent. Under within-dataset evaluation, balanced accuracy ranged from 67.8 to 87.3 on AD, from 48.5 to 81.3 on FTD, from 47.1 to 60.4 on FEP, and from 48.0 to 55.3 on MDD (Figure~\ref{fig:indomain_overview}). These results show that the downstream task placed a stronger constraint on performance than model family alone. This pattern is important because the evaluated self-supervised representations did not provide the same level of downstream separability across clinical tasks. Their usefulness should therefore be interpreted in the context of the specific clinical task and evaluated cohort rather than assumed to transfer uniformly across applications.


The strongest evidence for useful pretrained representations comes from AD and FTD. On AD, high performance was obtained by several model families, including simple classical baselines. The contribution of this task is therefore not evidence of superiority of one architecture, but evidence that the learned representations preserve information that supports strong participant level classification. FTD is more informative for comparing representations. Under LOSO, the proposed BrainLM-EEG variants were among the strongest models, with BrainLM-EEG-RoPE reaching 92.3 AUPRC and 92.9 balanced accuracy and BrainLM-EEG-Microstate reaching 94.6 and 92.3. This shows that the proposed representation can support strong out-of-sample classification on a task for which the model ranking was not already saturated across all approaches.

FEP and MDD produced a different pattern. Their results were weaker and less consistent, but they should not be treated as null findings. For FEP, several models produced AUPRC values above the task reference, although the preferred model depended on the protocol. For MDD, whose evaluation pool contained 26 patients and 72 controls, LOSO AUPRC values ranged from 30.7 to 42.5 relative to the 26.5 prevalence reference, with LaBraM at 42.5 and EEGNeX at 41.9. Balanced accuracy ranged from 49.5 to 61.4 against the nominal chance level of 50. The result is therefore not that the task contained no usable information, but that the available ranking signal translated less consistently into thresholded participant level classification.

The layer-wise analysis strengthens this interpretation without requiring a claim about the biological origin of the task differences. For MDD, the highest AUPRC was observed at layer 0 at 65.8, followed by 58.0, 51.0, 52.5, and 52.9 through the transformer. No transformer depth improved MDD AUPRC beyond the patch embedding representation. For AD and FTD the opposite happened, AUPRC kept climbing as the transformer added layers, reaching 95.5 and 92.8, respectively. Because each task has a different chance level, these raw numbers are not directly comparable, but the direction of change is, for MDD, the score is highest before the transformer runs and gets worse afterwards, while for AD and FTD it is lowest before the transformer runs and gets better afterwards. The relevant conclusion is therefore that the transformer increased downstream decodability much more clearly for AD and FTD than for MDD. The present experiments do not determine why the tasks differ, and the disease specific context for these differences has already been reviewed in Chapter~\ref{ch:background}.

\section{Evaluation Protocol and Participant Level Generalisation}\label{sec:d_protocol}

The comparison between five-fold cross-validation and LOSO provides an important methodological finding. The principal conclusion is not that one protocol reveals a uniquely correct model ranking, but that small clinical EEG cohorts can produce model rankings that depend strongly on how participants are resampled. The more reliable interpretation is therefore based on patterns that remain meaningful across evaluation settings rather than on small numerical differences within one split.

FTD provides the clearest example. Under five-fold cross-validation, BrainLM-EEG-RoPE reached 73.0 AUPRC and 48.5 balanced accuracy, below the nominal chance level of 50. Under LOSO, the same pretrained representation evaluated on the same 41 participants reached 92.3 and 92.9. BrainLM-EEG-Microstate similarly changed from 87.0 and 76.3 to 94.6 and 92.3 (Figure~\ref{fig:loso_vs_5fold}). With 41 participants, five-fold cross-validation results in both relatively small training sets and small test folds. The limited number of training participants can increase variability in the fitted model, while the small test folds make each individual participant influential on the fold level performance estimate. LOSO instead trains on nearly the full available cohort in each iteration, produces an out-of-sample prediction for every participant, and pools these predictions into a single participant level estimate.

The important finding is therefore twofold. First, the exact ordering of models on small cohorts should not be treated as stable when it changes with the resampling scheme. Second, the LOSO results still demonstrate that the proposed representations can generalise to held-out participants, particularly on AD and FTD. In this sense, the protocol comparison strengthens rather than weakens the thesis: it separates evidence that the representation contains useful participant level information from claims about small differences in model ranking.

\section{Generalisation Across Datasets}\label{sec:d_transfer}

The FEP-to-SCZ experiment extends the evaluation from unseen participants to an independently collected dataset. This distinction is central to the generalisation question. Good performance on held-out participants from the same dataset demonstrates participant level generalisation under a shared acquisition environment; external transfer additionally requires robustness to changes in cohort and acquisition conditions.

AUPRC decreased for every model when the FEP-trained models were evaluated on SCZ (Figure~\ref{fig:transfer_overview}). Using the SCZ prevalence reference of 39.0 given in Table~\ref{tab:chance_levels}, the highest transfer AUPRC under the five-fold setting was 54.4 for ATCNet and the highest under the all-data setting was 57.6 for EEGNeX. The proposed variants remained competitive in the first transfer setting, with values of 53.1 and 52.7, while their AUPRC values in the all-data setting lay between 42.2 and 47.4. Across the transfer table, balanced accuracy ranged from 47.4 to 63.4 against the nominal chance level of 50. The same challenge was observed for externally pretrained models: REVE, LaBraM, and CBraMod reached 46.1, 48.5, and 43.0 AUPRC, respectively, under the five-fold transfer protocol.

This result provides a clear boundary to the within-dataset findings. The proposed models learned representations that were useful for participant level classification, but those representations did not provide a consistent robustness advantage when both the cohort and acquisition setting changed. Importantly, this was not unique to BrainLM-EEG. Models pretrained on substantially larger EEG corpora showed the same difficulty, indicating that the FEP-to-SCZ shift remained challenging across modelling approaches.

The transfer experiment should nevertheless be interpreted in the context of the source task. The best within-dataset FEP AUPRC was 81.0, approximately 19 points above the FEP prevalence reference of 61.3. Because the source task was less strongly separated than AD or FTD, the experiment cannot isolate how much of the external performance reduction arose from the source task difficulty and how much arose from dataset shift. The conclusion that can be drawn directly is that, for this shift pair, participant level generalisation did not translate into a consistent cross-dataset robustness advantage. An external evaluation anchored on a stronger source task would provide a useful next test of this distinction.

\section{What the Representation Analyses Add}\label{sec:d_representation}

The analysis experiments provide evidence that goes beyond end-task model rankings. Together, the finetuning comparison, masking ablation, layer-wise analysis, CLS perturbation, and reconstruction experiment show how the pretrained representation can be used and where its useful information is located.

The finetuning comparison indicates that, with the labelled cohort sizes available here, the frozen representation was already a useful basis for downstream classification. Finetuning did not consistently improve on linear probing under the evaluated configuration (Figure~\ref{fig:lp_vs_ft}). This does not imply that finetuning is generally unsuitable. It does show that updating the full encoder was not necessary to obtain the strongest results observed here, and that the information learned during self-supervised pretraining could be accessed effectively using the lighter adaptation strategy. This is practically relevant for small clinical datasets, where unrestricted adaptation introduces many trainable parameters relative to the number of labelled participants.

The masking ablation provides the clearest evidence that the pretraining objective influences downstream representation quality. Masking whole channels rather than random spatiotemporal patches changed FTD AUPRC from 82.8 to 94.8 and changed AD AUPRC and balanced accuracy from 92.6 and 85.2 to 95.4 and 89.3 (Figure~\ref{fig:masking_ablation_results}). The FTD difference is 12 points, while the AD changes are approximately 3 to 4 points. The advantage was not universal, because FEP favoured random masking, and each pretraining configuration was evaluated from a single run. The conclusion is therefore not that one masking strategy is optimal for every EEG task, but that masking structure is a substantive design choice, not a neutral implementation detail. Similarly, the results with masking ratios of 75\% and the other evaluated settings show that the masking configuration inherited from vision models should be validated for EEG and not assumed to transfer unchanged.

The layer-wise experiment adds a second design insight. Downstream information did not simply accumulate until the final encoder output. The strongest AD result occurred at layer 3 and the strongest FTD result at layer 2, while the final representation was lower by 4 and 11 AUPRC points, respectively (Figure~\ref{fig:layerwise_results}). This means that selecting the representation depth is itself part of downstream adaptation. For clinical EEG tasks, the final transformer layer should not automatically be assumed to be the most informative representation.

The CLS perturbation provides direct evidence that the global token participates in reconstruction. Replacing it increased reconstruction error in all 71 held-out validation batches, by 3.86\% when zeroed and 4.42\% when replaced by Gaussian noise (Figure~\ref{fig:cls_perturbation_results}). The effect is modest relative to the roughly 71 visible patch tokens available to the decoder, so the CLS token should not be interpreted as the sole information bottleneck. Nevertheless, the consistent increase establishes that it carries information used by the decoder rather than acting as an unused downstream placeholder.

Finally, the reconstruction experiment shows that the learned encoder-decoder preserves information that remains useful to classifiers with very different inductive biases. AUPRC increased when LDA and ATCNet were trained on reconstructed rather than original EEG in all six reported comparisons, while balanced accuracy improved in only part of them (Figure~\ref{fig:reconstruction_asfeatures_results}). This is a particularly useful result because the effect is not tied to the downstream BrainLM classification head. It indicates that the reconstructed signal retains information relevant to participant ranking, as reflected by the higher AUPRC observed in the reported comparisons. The experiment does not establish a denoising mechanism, but it does provide independent evidence that the pretrained model learns a transformation that preserves clinically useful structure.

The architectural variants should be interpreted against these stronger representation level findings. BrainLM-EEG-Microstate and BrainLM-EEG-RoPE both produced strong results in individual settings, but neither modification showed a consistent improvement across tasks and protocols. This does not undermine the base representation learning approach. Instead, it indicates that the clearest gains in the present experiments came from choices such as masking and representation depth rather than from a universally superior positional or auxiliary loss modification. One possible explanation for the limited effect of the geometry-aware encoding is the common 19-electrode montage: with only 19 fixed spatial positions, a learned spatial embedding may already capture much of the available positional structure. This remains a hypothesis that should be tested with higher-density or variable-channel data.

\section{Implications for Self-Supervised EEG Modelling}\label{sec:d_implications}

Several practical conclusions follow from the combined experiments. First, evaluation should distinguish between generalisation to unseen participants and robustness to a new dataset. The former was demonstrated on the stronger downstream tasks, whereas the latter remained difficult in the FEP-to-SCZ experiment. Second, small clinical cohorts require care when model rankings are interpreted, because the FTD results show that rankings can change substantially with the resampling protocol. Third, the pretrained representation should not be treated as a fixed final layer object: masking strategy, representation depth, and adaptation method all affected how much useful information could be extracted.

These findings also clarify where self-supervised pretraining contributed most in the present study. Its value was strongest as a reusable representation learning stage. The frozen encoder supported competitive downstream performance, intermediate layers contained highly decodable information, and reconstructed EEG remained informative to independent classifiers. The results therefore support the use of self-supervised learning as a way of learning general EEG structure before labelled downstream adaptation, while showing that successful deployment still requires task specific choices about evaluation, representation extraction, and external validation.

\section{Limitations, Threats to Validity and Reproducibility}\label{sec:d_limitations}

The limitations affect the precision with which individual model differences can be interpreted, but they do not remove the broader patterns observed across the experiments. The first limitation is statistical. No significance testing was performed on the downstream model comparisons, and no confidence intervals were computed. Each pretraining configuration was also run with a single seed, so small differences between variants and masking configurations cannot be separated from run-to-run variation. The conclusions therefore focus on recurring patterns and larger directional changes rather than treating small numerical leads as evidence of general superiority.

The second limitation concerns the data. The downstream cohorts contain between 41 and 98 subjects, which is small relative to the number of parameters in the models. Only resting-state recordings were used, and all datasets were reduced to the common 19-channel montage. This standardisation supports comparison across datasets but sacrifices spatial resolution. Cross-dataset generalisation was assessed using a single external cohort, so the transfer conclusion is specific to one source-target pair.

The third limitation concerns the models and adaptation protocols. Finetuning used a single global learning rate with no layer-wise decay, no partial unfreezing, and no per-task search. Its lower performance relative to linear probing is therefore specific to that recipe. The classification head used for probing is a three-layer perceptron rather than a strictly linear classifier, so the relevant interpretation is downstream decodability, not linear separability. Comparisons with externally pretrained foundation models are also comparisons between complete pipelines, since those models differ in pretraining corpus, tokenisation, architecture, and input processing.

Several aspects of the experimental design strengthen the internal validity of the findings. All twelve models were evaluated on identical participant level partitions using the same aggregation and metric implementation. The partition between pretraining and downstream evaluation was made at the participant level, and the pretraining reserve mechanism described in Section~\ref{sec:m_datasets} prevented evaluation participants from entering the self-supervised pretraining corpus. These controls reduce the risk that the observed downstream performance is explained by participant level leakage.

\section{Answers to the Research Questions}\label{sec:d_rqs}

The four research questions stated in Section~\ref{sec:intro_rqs} can now be answered directly by considering the experiments together.

For the first research question, self-supervised pretraining produced useful downstream EEG representations. The proposed models were competitive with supervised, classical, and externally pretrained alternatives across several within-dataset evaluations and were among the strongest approaches on FTD under LOSO. They did not consistently outperform every baseline, so the evidence supports the usefulness and competitiveness of the learned representation rather than universal model superiority. Linear probing was the more effective adaptation strategy under the evaluated cohort sizes and hyperparameters.

For the second research question, the experiments demonstrate participant level generalisation but identify a clear distinction between within-dataset and cross-dataset robustness. LOSO produced out-of-sample predictions for every participant and showed strong performance on the best-performing tasks. In the FEP-to-SCZ evaluation, however, AUPRC decreased for every model relative to its within-dataset value, and no model family showed a consistent robustness advantage. Cross-dataset robustness therefore remains a more demanding problem than participant level generalisation within a shared dataset.

For the third research question, the experiments show that representation quality depends on pretraining and representation extraction choices. Masking ratio, masking strategy, and encoder depth all changed downstream performance. Intermediate layers produced the strongest decoding values for several tasks, and the canonical 75\% random masking setup was not uniformly best. The RoPE and microstate modifications did not produce a stable advantage across evaluation protocols, indicating that their benefit remains task dependent under the present setup.

For the fourth research question, the representation analyses provide direct evidence that BrainLM-EEG learns structured and reusable information. The CLS token contributes to reconstruction, task relevant information is distributed non-uniformly across encoder depth, and reconstructed EEG produced higher AUPRC in all six reported classifier comparisons. These findings show that the learned representation affects both latent space decodability and the information available after reconstruction.

Overall, the thesis demonstrates that self-supervised pretraining is a viable representation learning strategy for resting-state EEG and that its value is most clearly expressed through reusable participant level representations rather than through universal superiority of one architecture. The experiments identify concrete design choices that influence those representations, show that the learned information can support strong downstream classification in selected clinical tasks, and distinguish successful within-dataset generalisation from the more difficult problem of external dataset transfer. These conclusions provide a clear basis for improving future EEG foundation models: stronger external validation, task-aware masking, representation depth selection, and adaptation strategies matched to the size of the labelled cohort.
